We have built a laboratory, run a battery of controlled experiments, and
drawn an honest map. This is where we stand --- and where the
investigation goes next.

\section{What the experiment
established}\label{what-the-experiment-established}

At the core, the answer to the book's central question is now clear and
measured:

\begin{quote}
How much structured programming capability can a small language model
actually learn --- and how can we tell whether it learned rather than
memorized?
\end{quote}

The experiments separate cleanly into what the model \textbf{can} do and
what it \textbf{cannot}.

\textbf{It can learn, and it learns well --- within a narrow regime.} On
individual operations, the model reliably maps a natural-language spec
to correct Python, reaching \textbf{\textasciitilde0.90 pass@1}
in-distribution on a 0.88M parameter model. It generalizes across
\emph{instances}: on held-out prompts it passes, and on a canary it
produces the general program, not a memorized literal. So this is
genuine learning of the narrow task, not parroting.

\textbf{But it is surface-tied and non-compositional.} The two limits
that matter:

\begin{enumerate}
\def\labelenumi{\arabic{enumi}.}
\item
  \textbf{It anchors to phrasing, not meaning.} Change the variable
  names (\texttt{a,b} → \texttt{left,right}) and it drops to 0.35;
  change the wording (same computation, different words) and it falls to
  \textbf{0.000} --- below both baselines. The external reality check
  confirmed it: a human-written test set scored \textbf{0/16}, because
  the model learned a \emph{template}, not the \emph{task}.
\item
  \textbf{It cannot compose what it knows.} Train on \texttt{abs} and
  \texttt{add} separately, then ask for \texttt{abs\_add} --- it scores
  \textbf{0.000}. And the baselines also collapse to 0.000, so the split
  is fair: only genuine composition solves it, and the model has none.
\end{enumerate}

\textbf{The scaling conclusion is the sharpest negative result.}
Composition stays at \textbf{0.000} across model sizes from 0.88M to
29.3M. More capacity does not buy composition. This is not a capacity
ceiling in the \textasciitilde1--30M regime; it is a structural/learning
limit.

\section{The four-point grid
(measured)}\label{the-four-point-grid-measured}

\begin{longtable}[]{@{}lll@{}}
\toprule\noalign{}
Model & in-distribution & composition \\
\midrule\noalign{}
\endhead
\bottomrule\noalign{}
\endlastfoot
0.88M & 0.900 & \textbf{0.000} \\
6.5M & 0.900 & \textbf{0.000} \\
10.2M & 0.383 & \textbf{0.000} \\
29.3M & 0.897 & \textbf{0.000} \\
\end{longtable}

Four sizes, spanning \textasciitilde33× in capacity. Composition is
\textbf{0.000 at every one} --- a robust, falsifiable finding that the
boundary is not about how much brain the model has. Two honest notes on
the grid. First, the 10.2M point's \emph{in-distribution} regressed to
0.383 --- the deeper config was under-trained on this narrow task, an
instance of ``bigger isn't always better'' at tiny scale; it does not
change the composition result. Second, in-distribution performance is
\textbf{not monotonic} in capacity (0.900 → 0.900 → 0.383 → 0.897), a
worthwhile reminder that more capacity is not a strictly-increasing
guarantee even at these sizes.

\section{The methodology findings}\label{the-methodology-findings}

Independent of the model, the \textbf{methodology} produced its own
results --- and several are as valuable as any model number:

\begin{itemize}
\tightlist
\item
  \textbf{Test the tester.} Our own evaluator reported the model at
  0.000 on its first run. It was the \emph{judge} that was broken (a
  null-byte artifact), not the model. Fixing it revealed the true 0.897.
  Never trust a number from a system we have not audited.
\item
  \textbf{Test the teacher.} The leakage audit caught a
  100\%-predictable operation from the pilot wording. This is why
  surface splits need held-out template families.
\item
  \textbf{Test the data.} The diversity-vs-quantity experiment showed
  the \emph{same} example count produces surface-fragility (0.000) with
  minimal diversity or robust generalization (1.000) with high
  diversity. Diversity --- not quantity --- is what creates robust
  capability.
\item
  \textbf{Baselines are first-class.} A trivial TF-IDF retriever
  dominated the pilot task (1.000), proving it was ``too easy,'' and
  forced us to build harder splits. And the moments the model \emph{fell
  below} a baseline were the most informative results.
\item
  \textbf{The protocol, not capacity, determines continual learning.}
  Sequential training destroyed capability A (0.000) when learning B;
  replay preserved A perfectly (1.000) at no cost to B. This is a
  \emph{procedure} finding.
\end{itemize}

\section{The brewing experience,
distilled}\label{the-brewing-experience-distilled}

Beyond the numbers, there is a way of working this book was really
about. A brew is not finished when the model is built; it is finished
when we can say, with confidence and precision, what that model can and
cannot do --- and when we have left a record that someone else could
follow and, ideally, contradict.

That process had a rhythm. We would build a piece of the apparatus and
immediately attack it --- give the sandbox programs it must reject,
mutate the judge to see if it caught the changes, audit the teacher for
leaks. We would run an experiment, get a number, and then spend as much
effort asking \emph{how the model might have cheated} as we had spent
getting the number. And more than once, the interesting discovery was
not in the model's behavior but in our own tools: the judge that
mis-scored a working model, the data generator that leaked the
operation, the pilot task that a lookup table could solve perfectly.

This is the honest shape of empirical work. It is seldom a straight line
from hypothesis to clean result. It is a loop of building, testing the
builder, being surprised, and revising. The surprise is not a detour
from the method; it is the method. And a book about science should show
that loop, not hide it.

\section{Where this brew leads}\label{where-this-brew-leads}

The honest map points at specific next experiments, not vague ``more
data / bigger model'':

\begin{enumerate}
\def\labelenumi{\arabic{enumi}.}
\tightlist
\item
  \textbf{Does compositional exposure unlock composition?} The biggest
  open question. We never fed the model composited examples directly.
  Feeding a few (familiar) ones and testing unseen ones would tell us
  whether composition is learnable or a hard structure limit.
\item
  \textbf{Does the 10.2M in-distribution regression replicate?} The
  0.383 was unexpected. A proper training-duration and config sweep at
  10M would isolate whether it's under-training or a real depth penalty.
\item
  \textbf{Does continuing to scale past 30M shift the boundary?} The
  grid stops where the brew lab's CPU becomes impractical. A
  larger-scale run (or a GPU) would probe whether the composition wall
  eventually moves.
\item
  \textbf{Does diversity fix composition?} Diversity fixed the
  \emph{surface} problem (0.000 → 1.000). Whether the \emph{same} lever
  helps composition is the natural next test of the ``structural data''
  hypothesis.
\item
  \textbf{Human-style training.} The 0/16 human test set shows the model
  fails on naturally-phrased prompts. Training on diverse, human-style
  prompts --- the lever that already fixed surface robustness --- is the
  obvious path to real usability.
\end{enumerate}

\section{The honest ending}\label{the-honest-ending}

The model was the instrument. The experiment was the subject. We did not
build a general programmer, and we will not pretend otherwise. What we
\emph{did} build is the means to ask, precisely and honestly, what a
small model can and cannot learn --- and we got an answer worth having:

\begin{quote}
A small model learns narrow tasks well, but its knowledge is
surface-tied and non-compositional, and within the tested range neither
more capacity nor more data turns it into robust, compositional
capability.
\end{quote}

And the brakes we learned --- test the tester, test the teacher, compare
against baselines, audit shortcuts, preserve failures --- are the
durable part. That is the method worth keeping, and it's the reason the
book's real subject was never the model. The brew was never only about
the model; it was about learning to ask a question carefully enough that
the answer could be trusted.
